%=========================================
% Introduction
%=========================================

\section{Introduction}

\subsection{Background}

Headache disorders are common health conditions that can substantially affect
quality of life, daily functioning, and healthcare utilization. In clinical
studies, headache burden is often measured using count outcomes, such as the
number of days on which a patient experiences a headache during a specified
observation period.

The Poisson distribution is commonly used to model count outcomes because it
provides a natural probabilistic framework for the occurrence of events. A
conventional Poisson model, however, assumes that the conditional mean and
variance of the outcome are equal. This assumption may be inappropriate when
patients differ substantially in their underlying susceptibility to headaches.

When observations arise from a population containing individuals with
different underlying headache frequencies, the resulting data may exhibit
latent heterogeneity. One consequence of such heterogeneity can be
over-dispersion, whereby the observed variance exceeds the mean. Ignoring
this structure and applying a single-population Poisson model may therefore
result in an inadequate representation of the observed count distribution.

Finite mixture models provide a flexible framework for modelling such latent
heterogeneity. Rather than assuming that all observations arise from a single
population, a finite mixture model assumes that the population consists of a
finite number of unobserved subpopulations, or mixture components. A Poisson
finite mixture model (PFMM) allows each component to have its own Poisson rate
while simultaneously estimating the proportion of individuals belonging to
each latent component.

This study applies Poisson finite mixture modelling to headache frequency data
from a randomized clinical trial. The analysis considers the distribution of
headache frequency, evaluates evidence of over-dispersion, estimates Poisson
finite mixture models, and identifies an appropriate number of latent
components. Nonparametric maximum likelihood estimation (NPMLE) and the
expectation-maximization (EM) algorithm are used in the estimation process,
while posterior probabilities are used to classify patients into the
identified latent subpopulations.

Patient age is additionally considered as a covariate to investigate whether
an observed patient characteristic can explain differences in headache
frequency across the latent mixture components. The analysis is implemented
using both R and SAS, providing a reproducible computational framework for
the statistical investigation.

\subsection{Problem Statement}

Headache frequency may vary considerably between patients, even when the
patients are enrolled in the same clinical study. Such variation may reflect
underlying differences in susceptibility to headaches and may result in a
population that is not adequately described by a single probability
distribution.

A standard Poisson model assumes that all observations arise from one
population with a common underlying event rate, subject to the model's
covariate structure. When the population contains unobserved subgroups with
different event rates, this assumption may be violated and the resulting data
may exhibit over-dispersion.

The statistical problem addressed in this study is therefore to determine
whether the observed heterogeneity in headache frequency can be represented
by a finite mixture of Poisson distributions and to identify the number of
latent components that provides an appropriate representation of the data.

A related objective is to investigate whether patient age helps explain the
differences in headache frequency associated with the identified mixture
components. This provides an opportunity to distinguish heterogeneity that
can be explained by an observed covariate from heterogeneity that may remain
unexplained by the available data.

\subsection{Research Objectives}

\subsubsection{General Objective}

The general objective of this study is to investigate latent heterogeneity in
headache frequency using Poisson finite mixture models and to evaluate their
ability to represent over-dispersed count data.

\subsubsection{Specific Objectives}

The specific objectives are to:

\begin{enumerate}[label=\arabic*.]
    \item Describe the distribution and descriptive characteristics of
    headache frequency in the study population.

    \item Assess whether the observed headache frequency data exhibit
    over-dispersion relative to a conventional Poisson model.

    \item Estimate Poisson finite mixture models with different numbers of
    latent components.

    \item Compare the fitted mixture models and identify an appropriate
    number of latent subpopulations.

    \item Estimate the parameters of the selected mixture model using
    nonparametric maximum likelihood estimation and the
    expectation-maximization algorithm.

    \item Classify patients into the identified latent subpopulations using
    posterior probabilities.

    \item Investigate whether patient age explains differences in headache
    frequency across the identified mixture components.

    \item Implement the statistical analysis using both R and SAS.
\end{enumerate}

\subsection{Research Questions}

The study addresses the following research questions:

\begin{enumerate}[label=\arabic*.]
    \item What is the distribution of headache frequency among the patients
    in the study population?

    \item Does the headache frequency data exhibit over-dispersion relative
    to a single-population Poisson model?

    \item How many latent subpopulations provide an appropriate
    representation of the observed headache frequency distribution?

    \item What are the estimated Poisson rates and mixing proportions of the
    selected finite mixture model?

    \item How can patients be classified into the identified latent
    subpopulations using posterior probabilities?

    \item Does patient age help explain differences in headache frequency
    across the latent mixture components?
\end{enumerate}

\subsection{Scope of the Study}

This study focuses on the statistical modelling of headache frequency
recorded during a four-week baseline period. The primary outcome is the
number of headache days, treated as a count variable.

The analysis focuses on Poisson-based finite mixture modelling, including
assessment of over-dispersion, estimation and comparison of mixture models,
latent subpopulation classification, and the incorporation of patient age as
a covariate.

The study is concerned with statistical description and inference rather than
the evaluation of treatment effectiveness. Consequently, the analysis does
not attempt to establish causal relationships between patient characteristics
and headache frequency.

The computational analysis is conducted using R and SAS, while \LaTeX{} is
used to prepare and document the research report.

%============================================================
\subsection{Structure of the Report}
%============================================================

The remainder of this report is organized as follows.

Section 2 describes the study data and presents the exploratory and descriptive
analyses.

Section 3 presents the literature review, covering count-data modelling,
heterogeneity and finite mixture models, nonparametric maximum likelihood
estimation, estimation of finite mixture models, and approaches to determining
the number of components.

Section 4 presents the statistical methodology, including the Poisson
distribution, assessment of over-dispersion, Poisson finite mixture models,
nonparametric maximum likelihood estimation, the expectation-maximization
algorithm, and posterior probability-based classification.

Section 5 presents the model estimation, model comparison, and principal
statistical results.

Section 6 discusses the findings in relation to the research objectives,
including the role of patient age and the implications of the identified
latent heterogeneity.

Section 7 summarizes the main findings and presents the conclusions and
implications of the study.